\documentclass[12pt]{article}
\usepackage[utf8]{inputenc}
\usepackage{amsmath}
\usepackage{amsfonts}
\usepackage{geometry}
\usepackage{titlesec}

\geometry{a4paper, margin=1in}

\title{Especificación Técnica: Modelo de Matriz de Conceptos Recursivos (ConceptMatrix 3D)}
\author{Arquitectura de Redes Neuronales Dispersas}
\date{Marzo 2026}

\begin{document}

\maketitle

\section{Introducción}
El modelo propuesto se define como una estructura de datos tridimensional dispersa, denotada como $CM$, donde cada coordenada activa $C(x, y, z)$ no solo representa un identificador semántico, sino que actúa como el punto de entrada a una red neuronal local $NN_c$.

\section{Estructura Espacial}
La matriz se define en un espacio euclidiano de tres dimensiones con una resolución $S$:
$$CM \in \mathbb{R}^{S \times S \times S}$$
Dada la naturaleza dispersa de la matriz, el conjunto de conceptos definidos $K$ es un subconjunto tal que $|K| \ll S^3$.

\section{El Nodo Conceptual como Red Neuronal}
Cada concepto $k \in K$ posee una función de transformación interna:
$$f_k(\mathbf{v}_{in}) = \sigma(W_k \cdot \mathbf{v}_{in} + b_k)$$
Donde:
\begin{itemize}
    \item $\mathbf{v}_{in}$ es el vector de contexto proveniente de conceptos adyacentes o de la entrada del sistema.
    \item $W_k$ es la matriz de pesos específica del concepto $k$, que representa miles de neuronas locales.
    \item $\sigma$ es la función de activación (ej. ReLU) que determina la relevancia del concepto en el flujo actual.
\end{itemize}

\section{Recursividad y Definición de Conceptos}
A diferencia de los modelos tradicionales, la salida de un nodo $C_i$ es una lista de punteros a otras coordenadas dentro de la misma matriz $CM$. La definición de un concepto "Carro" se expresa como:
$$\text{Def}(C_{carro}) = \{C_{motor}, C_{rueda}, C_{chasis}, \dots, C_n\}$$
Donde cada $C_i$ es, a su vez, una instancia de una red neuronal independiente.

\section{Proceso de Entrenamiento}
El entrenamiento del modelo se basa en la minimización de la distancia entre la trayectoria predicha por las redes neuronales y la trayectoria real definida en el corpus:

\subsection{Función de Pérdida (Loss Function)}
Para una secuencia de conceptos, la pérdida se calcula mediante la distancia euclidiana en el espacio 3D:
$$\mathcal{L} = \sum_{t=1}^{T} \| \hat{P}_t(x,y,z) - P_t(x,y,z) \|^2$$
Donde $\hat{P}_t$ es la coordenada predicha por la red del concepto actual y $P_t$ es la coordenada real del siguiente concepto en la definición.

\section{Conclusión}
Este modelo permite una navegación semántica basada en lógica estructural. La densidad de neuronas por cada punto de la matriz garantiza que el sistema no solo almacene posiciones, sino comportamientos dinámicos ante diferentes contextos lingüísticos.

\section{Arquitectura Interna del Nodo: Micro-Capas de Procesamiento}
Cada punto activo en la matriz $CM$ no es una unidad atómica, sino un contenedor de una arquitectura neuronal densa. Si definimos $N$ como el número de neuronas por concepto (donde $N \approx 1000$), el estado interno de un concepto se describe como un tensor de parámetros locales.

\subsection{Procesamiento de Contexto Local}
Cuando un concepto $C_k$ es activado, sus $N$ neuronas realizan una operación de filtrado selectivo sobre las señales de entrada $\mathbf{v}_{in}$. Esto permite que el concepto responda de manera distinta según el contexto:
$$ \mathbf{h}_k = \phi( \mathbf{W}_{local}^{(k)} \cdot \mathbf{v}_{in} + \mathbf{b}_{local}^{(k)} ) $$
Donde $\mathbf{h}_k \in \mathbb{R}^N$ representa el vector de estado oculto del concepto.

\subsection{Selección de Definición (Attention Mechanism)}
Dado que la definición de un concepto "Carro" contiene múltiples sub-conceptos $\{\text{motor, rueda, ...}\}$, las mil neuronas internas actúan como un mecanismo de atención que asigna pesos de relevancia $\alpha$ a cada hijo de la definición:
$$ \alpha_i = \text{softmax}( \text{score}(\mathbf{h}_k, C_i) ) $$
Esto permite que el modelo, ante la entrada "ruido en el motor", active con mayor intensidad la red neuronal del sub-concepto $C_{motor}$ que la del sub-concepto $C_{chasis}$.

\subsection{Aprendizaje de Micro-Pesos}
Durante la etapa de entrenamiento, el gradiente del error $\nabla \mathcal{L}$ se distribuye no solo para mover la posición $(x, y, z)$ del concepto en la matriz dispersa, sino para ajustar los miles de pesos internos $\mathbf{W}_{local}^{(k)}$. Esto permite que el concepto "aprenda" su propia semántica interna independientemente del resto de la matriz.

\section{Resolución de Limitaciones Intrínsecas de los LLMs Tradicionales}

La arquitectura de Matriz de Conceptos Recursiva (RCM) ofrece soluciones estructurales a los problemas de los modelos transformadores convencionales:

\subsection{Mitigación de Alucinaciones mediante Anclaje Semántico}
En los LLMs actuales, el conocimiento está distribuido de forma difusa. En este modelo, el conocimiento está \textit{localizado} en coordenadas $(x, y, z)$ y validado por su lista de sub-conceptos.
\begin{itemize}
    \item \textbf{LLM Tradicional:} Genera texto basado en la secuencia más probable, aunque sea falsa.
    \item \textbf{RCM:} Si la red neuronal del concepto "Carro" no tiene una conexión activa hacia el concepto "Volar", el modelo tiene una restricción física/lógica que le impide generar esa afirmación, eliminando la alucinación por diseño estructural.
\end{itemize}

\subsection{Razonamiento Simbólico-Neuronal (Neuro-Symbolic)}
Los LLMs tienen dificultades con la lógica multietapa. Al utilizar una estructura recursiva, este modelo realiza un "recorrido de grafo":
$$ \text{Pensamiento} = \text{Trayectoria}(C_{inicio} \to C_{def1} \to C_{sub-def} \dots) $$
Cada "salto" entre coordenadas requiere una activación de las mil neuronas locales, lo que obliga al modelo a mantener una coherencia lógica paso a paso, similar al razonamiento humano sistemático.

\subsection{Eficiencia Computacional y Actualización de Conocimiento}
Para enseñar un nuevo dato a un LLM actual, se debe re-entrenar todo el modelo (o gran parte de él). En la arquitectura RCM:
\begin{itemize}
    \item \textbf{Aprendizaje Localizado:} Si se descubre una nueva característica del "Motor", solo se actualizan los pesos $W_{local}^{(motor)}$ y su vector de coordenadas. El resto de la matriz permanece intacto.
    \item \textbf{Consumo de Recursos:} Solo se activan los conceptos relevantes para la consulta, en lugar de procesar billones de parámetros para cada palabra.
\end{itemize}

\end{document}